\documentclass[conference]{IEEEtran}
\IEEEoverridecommandlockouts

\usepackage{cite}
\usepackage{amsmath,amssymb,amsfonts}
\usepackage{algorithmic}
\usepackage{graphicx}
\usepackage{textcomp}
\usepackage{xcolor}
\usepackage{booktabs}
\usepackage{url}
\usepackage{hyperref}

\def\BibTeX{{\rm B\kern-.05em{\sc i\kern-.025em b}\kern-.08em
    T\kern-.1667em\lower.7ex\hbox{E}\kern-.125emX}}

\begin{document}

\title{Regime-Adaptive Multimodal Transformer for
Cross-Market Equity Forecasting}

\author{\IEEEauthorblockN{Vivek (230119)}
\IEEEauthorblockA{\textit{B.Tech in Computer Science and Artificial Intelligence} \\
\textit{Rishihood University}\\
Sonipat, India \\
vivek.23csai@nst.rishihood.edu.in}
\and
\IEEEauthorblockN{Shivansh (230054)}
\IEEEauthorblockA{\textit{B.Tech in Computer Science and Artificial Intelligence} \\
\textit{Rishihood University}\\
Sonipat, India \\
shivansh.g23csai@nst.rishihood.edu.in}
}
\maketitle

\begin{abstract}
We test whether a regime-adaptive multimodal transformer (RAMT) can rank NIFTY~200
constituents cross-sectionally on a 21-day horizon. RAMT fails: predictions collapse
($\sigma \approx 0.0015$ across names), validation IC turns negative, and a
60-line LightGBM diagnostic on the same features produces $\mathrm{IC}=+0.021$.
We replace RAMT with two alternatives: (a)~plain 21-day momentum gated by a 3-state
Hidden Markov Model on the index, and (b)~Chronos-T5-Small fine-tuned with LoRA
($r{=}8$) on the same 10 features. Across the 2024--2026 out-of-sample window,
Foundation-Only Chronos-LoRA delivers Sharpe $1.34$, the strongest single component;
Momentum$+$HMM delivers Sharpe $0.83$ with a more conservative drawdown profile;
the regime-gated Triple-Expert hybrid \emph{underperforms} both at Sharpe $0.54$.
Across four historical HMM ablation windows (2008--2010, 2010--2012, 2013--2015,
2024--2026) the regime gate saves $9.4$pp of drawdown in the 2008 crisis and
turns a Sharpe of $-3.0$ into $+0.79$ in 2010--2012 but caps upside in clean bulls.
Our three findings: (1)~simpler architectures beat deeper ones at sub-million-parameter
scales on this universe; (2)~the RAMT failure is architectural, not a feature
problem (LightGBM and momentum found signal on the same panel); (3)~the HMM is
\emph{conditional insurance}, not always-on alpha.
\end{abstract}

\begin{IEEEkeywords}
quantitative finance, factor investing, hidden markov models, foundation models,
LoRA, Chronos-T5, regime detection, NIFTY 200, ablation
\end{IEEEkeywords}

% --------------------------------------------------------------------
\section{Introduction}
% --------------------------------------------------------------------
The original hypothesis of this project was that a multimodal transformer with
explicit regime cross-attention and a tournament ranking loss would extract
cross-sectional alpha from NIFTY~200 names that simpler factor models miss. The
hypothesis is rejected by the data. After two attempts at training the
transformer (RAMT) we observed prediction-variance collapse and negative
validation IC. A diagnostic with plain LightGBM on the same ten features
produced statistically meaningful information coefficient, ruling out feature
adequacy as the cause. The remainder of the project tests two simpler
alternatives and a foundation-model fine-tune.

The contribution of this paper is empirical: we report a complete ablation
across one ML-only baseline (Momentum$+$HMM), one transformer-from-scratch
attempt (RAMT), one foundation-model fine-tune (Chronos-T5-Small with LoRA
adapters), and two interaction modes between them, on the same 2024--2026
out-of-sample window with identical friction and risk-control rules. We then
report a separate four-window study isolating HMM regime sizing against flat
sizing.

% --------------------------------------------------------------------
\section{Related Work}
% --------------------------------------------------------------------
Cross-sectional momentum on equities was established by Jegadeesh and
Titman~\cite{jegadeesh1993} and refined into the global ``everywhere''
formulation by Asness, Moskowitz, and Pedersen~\cite{asness2013}. Hidden Markov
Models for regime detection follow Hamilton's regime-switching framework~\cite{hamilton1989},
extended to multivariate financial settings by Guidolin and
Timmermann~\cite{guidolin2007}. Transformer architectures originate from
Vaswani et al.~\cite{vaswani2017} and were specialized for time-series
forecasting in Lim et al.~\cite{lim2021}'s Temporal Fusion Transformer. Our
Phase~3 substitutes a pre-trained foundation model: Chronos~\cite{chronos2023}
treats time-series tokenization with the T5 architecture, and we adapt it
parameter-efficiently using LoRA~\cite{lora2021}, which constrains updates
to low-rank decompositions $W = W_0 + BA$ with $B \in \mathbb{R}^{d \times r}$,
$A \in \mathbb{R}^{r \times k}$, $r \ll \min(d,k)$.

% --------------------------------------------------------------------
\section{Data and Features}
% --------------------------------------------------------------------
\textbf{Universe.} A static 2026 NIFTY~200 snapshot. We acknowledge survivorship
and constituent-drift bias as a limitation rather than a fix.

\textbf{Sources.} Daily OHLCV per ticker via \texttt{yfinance}; index series for
\texttt{\^{}NSEI} for the regime feed and benchmark; sector mapping is a hand-curated
table in \texttt{features/sectors.py}.

\textbf{Features (10).} Per ticker per date, lagged at $T{-}1$ to enforce
information frontiers: 21-day return (\texttt{Ret\_21d}), 5-day return,
realized volatility (21-day), volume z-score, price relative to 50-day SMA,
price relative to 200-day SMA, 14-day RSI, MACD histogram, NIFTY 21-day return
(macro feed), NIFTY realized volatility (macro feed). The label is 21-day
forward sector-neutral alpha (\texttt{Sector\_Alpha}).

\textbf{Modalities (4).} Price, technicals, volume, macro. The HMM is fit on the
NIFTY index alone (3-state Gaussian, fit by Baum--Welch).

\textbf{Manifest.} \texttt{data/manifest.csv} ships 201 parquet entries with
md5 hashes; \texttt{scripts/check\_pipeline\_health.py --verify-manifest}
re-hashes them on demand for reproducibility.

% --------------------------------------------------------------------
\section{Phase~1: Baseline Models}
% --------------------------------------------------------------------
We first tested daily-return prediction with XGBoost and an LSTM. Both failed:
directional accuracy 49.4\% (XGBoost) and 48.9\% (LSTM) on the walk-forward
test set; mean information coefficient $-0.061$ and $-0.041$ respectively
(\texttt{results/backtesting/phase1\_baselines/}). The label was the bug, not
the algorithm: daily returns at this granularity are dominated by noise. We
pivoted to 21-day forward sector-neutral alpha for Phases~2 and~3.

% --------------------------------------------------------------------
\section{Phase~2: RAMT Architecture and Failure}
% --------------------------------------------------------------------
\subsection{Architecture}
RAMT comprises six modules:
(1)~a \texttt{MultimodalEncoder} for the four modalities;
(2)~standard sinusoidal positional encoding;
(3)~\texttt{RegimeCrossAttention}, where the regime embedding is the query;
(4)~a 3-expert MoE block (bull, high-vol, bear flavor) followed by transformer layers;
(5)~dual prediction heads for monthly and daily horizons; and
(6)~a tournament ranking loss with an auxiliary daily MSE anchor.

\subsection{Training Setup}
Walk-forward, 6-month refresh, sequence length 30. After speed and
overfitting fixes (workers $=4$, cache $=250$, MoE collapsed to a single
expert), the final architecture has 131{,}298 parameters across 2~transformer
layers and 8~heads.

\subsection{Failure Analysis}
Across 25~epochs and four walk-forward folds, validation correlation to
\texttt{Sector\_Alpha} went \emph{negative} ($-0.036$) and worsened
monotonically. Per-date prediction standard deviation across the 200 names
collapsed to $\sigma \approx 0.0015$. The mechanism: the tournament ranking
loss has a fixed-point at zero scores --- once predictions collapse, pairwise
margins vanish, gradients vanish, and the model has no signal to escape. We
later raised the daily-MSE anchor weight from 0.05 to 0.20, restored capacity
to 131k parameters, and re-ran 7~epochs; the negative-correlation curve was
already monotonic and we stopped.

Final RAMT metrics on the blind 2024--2026 window:
$\mathrm{DA}=47.2\%$, $\mathrm{RMSE}=0.072$, $\mathrm{mean\ IC}=-0.016$,
backtested $\mathrm{Sharpe}=0.485$, $\mathrm{MaxDD}=-6.4\%$
(\texttt{results/models/ramt/ramt\_metrics.json}).

\subsection{Diagnostic that motivated the pivot}
A 60-line LightGBM script on the same 10 features and the same train/test
split produced $\mathrm{IC}=+0.021$, $\mathrm{IR}=+6.41$, top-5 positive-alpha
rate $53.8\%$. Pure 21-day momentum without any model produced
$\mathrm{IC}=+0.0065$, top-5 spread $+1.59\%$ per month, top-5 positive rate
$56.3\%$. Both beat RAMT. The features carried signal; the transformer did not
extract it.

% --------------------------------------------------------------------
\section{Phase~3: Foundation Model with LoRA}
% --------------------------------------------------------------------
\subsection{Architecture}
We freeze the pre-trained Chronos-T5-Small backbone (${\sim}190\mathrm{M}$
parameters) and attach LoRA adapters with rank $r{=}8$ to the attention
projection matrices, training only the adapters and a small ranking head ---
together approximately $0.1\%$ of the backbone's parameter count. Input is the
same $(B, 30, 10)$ feature tensor used by RAMT. Output is a scalar alpha score
per ticker per rebalance date.

\subsection{Training Setup}
LoRA $r{=}8$, $\alpha{=}16$, dropout $0.1$. Optimizer AdamW, learning rate
$1\mathrm{e}{-4}$, batch size 64, walk-forward training. Reference scripts:
\texttt{models/lora\_experiment/chronos\_lora\_v2.py},
\texttt{scripts/train\_chronos\_lora.py}. The trained adapter is committed to
\texttt{results/models/lora/chronos\_lora\_adapter.pt}.

\subsection{Results}
Out-of-sample directional accuracy $\mathrm{DA}=47.5\%$ (LoRA v1) /
$47.0\%$ (LoRA v2), but mean IC turns positive
($+0.0090$ / $+0.0022$) --- a sign change relative to RAMT. More importantly,
prediction standard deviation no longer collapses
($\sigma_{\mathrm{pred}} \approx 4.7\mathrm{e}{-3}$ for LoRA versus
${\sim}1.5\mathrm{e}{-3}$ for RAMT, a ${\sim}3\times$ widening that allows
ranking to actually separate names). Backtested Sharpe on 2024--2026 reaches
$1.34$ (Foundation-Only), the strongest single-component result in the
project. See Table~\ref{tab:ablation}.

% --------------------------------------------------------------------
\section{Final Strategy: Momentum + HMM + Sector Cap}
% --------------------------------------------------------------------
The currently-deployed strategy is intentionally simpler than the ablation
winner. Each rebalance date (every 21 trading days):
\begin{enumerate}
\item Sort the 200 names by 21-day return.
\item Take top 5 with at most one name per NSE sector.
\item Size the book by the current HMM regime: 100\% / 50\% / 20\%
      for Bull / High-Vol / Bear.
\item Apply per-stock 7\% stop and a 15\% portfolio drawdown killswitch.
\item Charge friction at 0.22\% of traded notional.
\end{enumerate}
Out-of-sample (\texttt{2024-01-10} to \texttt{2026-01-27}, 2-year window):
$\mathrm{Sharpe} = 0.83$, $\mathrm{CAGR} = 13.5\%$,
$\mathrm{MaxDD} = -18.7\%$, win rate $64\%$. Versus NIFTY buy-and-hold over
the same calendar (Sharpe $0.65$, CAGR $7.7\%$, MaxDD $-15.8\%$): $+0.18$
Sharpe, $+5.8$pp CAGR, $-2.9$pp on max drawdown.

% --------------------------------------------------------------------
\section{Ablation}
% --------------------------------------------------------------------

\subsection{Cross-component ablation (2024--2026 window)}
\begin{table}[ht]
\centering
\caption{Net-of-friction performance on the 2024--2026 out-of-sample window.}
\label{tab:ablation}
\begin{tabular}{lccc}
\toprule
\textbf{Variant} & \textbf{Sharpe} & \textbf{CAGR} & \textbf{MaxDD} \\
\midrule
Momentum + HMM (production)            & 0.83 & 13.5\% & -18.7\% \\
RAMT transformer                       & 0.49 & n/a    & -6.4\%  \\
Chronos-LoRA Foundation-Only           & \textbf{1.34} & 23.5\% & -16.0\% \\
Simple Hybrid (50/50 Mom + Chronos)    & 0.91 & 22.8\% & -11.1\% \\
Triple-Expert (Mom + Chronos + HMM)    & 0.54 & 8.2\%  & -13.8\% \\
\midrule
\multicolumn{4}{l}{\emph{Component-removal cells:}} \\
Hybrid $-$ HMM (alias of Simple Hybrid)    & 0.91 & 22.8\% & -11.1\% \\
Hybrid $-$ Chronos (alias of Mom+HMM prod) & 0.83 & 13.5\% & -18.7\% \\
Hybrid $-$ Momentum                        & \multicolumn{3}{c}{not run --- module missing} \\
\bottomrule
\end{tabular}
\end{table}

The two interpretable component-removal cells are aliases of existing
scenarios because, by construction, removing HMM from the Triple-Expert
collapses to the 50/50 hybrid, and removing Chronos collapses to
Momentum$+$HMM. The third cell (Chronos$+$HMM, no momentum) is honestly
recorded as \emph{not run}: the in-repo \texttt{models/hybrid\_backtester.py}
module referenced from \texttt{scripts/run\_diagnostic\_ablation.py} is absent
in the current working tree, and we chose not to fabricate a number for the
table.

\subsection{HMM regime-sizing ablation across four historical windows}

\begin{table}[ht]
\centering
\caption{HMM-conditioned vs flat sizing, four windows.}
\label{tab:hmm}
\begin{tabular}{lcccc}
\toprule
\textbf{Window} & \textbf{HMM Sharpe} & \textbf{Flat Sharpe} & \textbf{HMM DD} & \textbf{Flat DD} \\
\midrule
2008--2010 (GFC)   & 0.25 & 0.17  & -43.2\% & -52.7\% \\
2010--2012         & 0.79 & -3.00 & -12.2\% & -30.1\% \\
2013--2015         & 0.84 & 0.59  & -33.2\% & -33.2\% \\
2024--2026 (bull)  & 0.66 & 1.35  & -18.7\% & -19.7\% \\
\bottomrule
\end{tabular}
\end{table}

The pattern is consistent across four windows: HMM gating helps in volatile
and crisis windows (saves 9.4pp drawdown in 2008, rescues a $-3.0$ Sharpe in
2010--2012) and \emph{hurts} in clean bull markets (cuts Sharpe from $1.35$ to
$0.66$ in 2024--2026). Treating HMM as conditional insurance rather than
always-on alpha is the central architectural lesson of the project.

% --------------------------------------------------------------------
\section{Limitations}
% --------------------------------------------------------------------
\textbf{Survivorship bias.} The universe is a 2026 NIFTY~200 snapshot, not the
historically reconstructed index membership. Although the panel is
time-aligned, the cross-section excludes names that were in the index but
later delisted.

\textbf{Window length.} The headline 2024--2026 Phase~3 numbers are from a
single two-year out-of-sample window. The strategy has not been tested
through a 2008-style or 2020-style crash in this configuration; the
four-window HMM study is the closest stress test.

\textbf{Stop-loss is a no-op in places.} The 7\% per-stock stop and 15\%
portfolio killswitch were rarely triggered in the 2024--2026 window, so the
backtest does not strongly exercise that path.

\textbf{Friction model.} Costs are modeled as a fixed 0.22\% of traded
notional per rebalance. Real Indian small/mid-cap impact is often higher and
varies with name and size; this is a known optimistic bound.

\textbf{Static feature set.} We did not search over feature subsets or
include text/sentiment features in the final ablation, although a FinBERT
sentiment pathway exists in the repository.

\textbf{Incomplete component-removal ablation.}
The Hybrid $-$ Momentum ablation cell, which would isolate the contribution
of Chronos-LoRA predictions without the momentum ranking signal, could not
be computed at submission time due to a missing backtesting module
dependency. This cell is recorded as \texttt{not\_run} in the ablation
summary JSON and represents an open validation gap. All other
component-removal cells are populated from real experimental artifacts.

% --------------------------------------------------------------------
\section{Conclusion}
% --------------------------------------------------------------------
Three findings:

\emph{Simpler beat deeper.} At our data and parameter scale, plain 21-day
momentum outperformed a 131k-parameter regime-aware transformer trained
end-to-end. Foundation-Only Chronos-LoRA, which leans on a 190M-parameter
pre-trained backbone with only ${\sim}0.1\%$ trainable, then beat both. The
inflection happened at the foundation-model boundary, not by adding
architectural complexity to the from-scratch model.

\emph{The RAMT failure was architectural.} The same 10 features driven into
LightGBM produced positive IC; driven into 21-day momentum produced positive
top-5 spread. The transformer's tournament ranking loss collapsed to a
fixed-point and the gradients went to zero. We document this honestly rather
than masking it with longer training.

\emph{HMM is conditional insurance.} Across four historical windows the HMM
regime gate adds Sharpe and reduces drawdown in crisis and high-volatility
windows, and subtracts Sharpe in clean bulls. The honest deployment policy is
to gate \emph{only} during volatile regimes; using HMM as an always-on
multiplier is what produced the Triple-Expert Sharpe of 0.54.

The natural next step is a learned regime-conditional gate that defers to
Foundation-Only in calm regimes and to HMM-protected sizing in volatile ones,
with the gate trained on regime-conditional out-of-sample Sharpe rather than
hand-tuned weights.

% --------------------------------------------------------------------
\bibliographystyle{IEEEtran}
\bibliography{references}

\end{document}
